\documentclass[theme=teal,margin=normal,11pt]{khmermath}

\school{វិទ្យាល័យ ព្រះស៊ីសុវត្ថិ}
\teacher{ហេង}
\phone{077 453 282}
\subject{គណិតវិទ្យា ថ្នាក់ទី១២}

\enablebanner

\begin{document}

	% ចំណងជើងធំ
	\lessontitle[វិញ្ញាសាត្រៀមប្រឡងសញ្ញាបត្រមធ្យមសិក្សាទុតិយភូមិ]{វិញ្ញាសាគណិតវិទ្យាពិសេស}

	\begin{formulabox}{រូបមន្តគន្លឹះត្រូវចាំ (Key Formulas)}
		\begin{multicols}{2}
			\begin{itemize}[leftmargin=15pt, itemsep=2pt]
				\item $\lim\limits_{x \to 0} \frac{\sin x}{x} = 1$ \quad និង \quad $\lim\limits_{x \to 0} \frac{1 - \cos x}{x^2} = \frac{1}{2}$
				\item $(uv)' = u'v + uv'$ \quad និង \quad $\left(\frac{u}{v}\right)' = \frac{u'v - uv'}{v^2}$
				\item $\int x^n \, dx = \frac{x^{n+1}}{n+1} + c \quad (n \neq -1)$
				\item $\int e^{ax} \, dx = \frac{1}{a} e^{ax} + c$
			\end{itemize}
		\end{multicols}
	\end{formulabox}

	\vspace{4pt}

	\begin{sectionbox}{ប្រធានលំហាត់ (Exam Problems)}
		\begin{circleenum}
			\item \textbf{(២.៥ ពិន្ទុ)} គណនាលីមីតខាងក្រោម៖
			\begin{multicols}{2}
				\begin{khmerenum}
					\item $A = \lim\limits_{x \to 2} \frac{x^3 - 8}{x^2 - 4}$
					\item $B = \lim\limits_{x \to 0} \frac{\sqrt{1+3x} - 1}{x}$
					\item $C = \lim\limits_{x \to +\infty} \left(\sqrt{x^2+4x} - x\right)$
					\item $D = \lim\limits_{x \to 0} \frac{1 - \cos 2x}{x \sin x}$
				\end{khmerenum}
			\end{multicols}

			\vspace{3pt}
			\item \textbf{(២.៥ ពិន្ទុ)} សិក្សាភាពជាប់នៃអនុគមន៍ $f(x)$ ត្រង់ $x = 0$ ៖
			\[
			f(x) = \begin{cases}
				\frac{\sin 3x}{x} & \text{បើ } x > 0 \\
				3 & \text{បើ } x = 0 \\
				x^2 + 3 & \text{បើ } x < 0
			\end{cases}
			\]

			\vspace{3pt}
			\item \textbf{(៥.០ ពិន្ទុ)} គេឱ្យអនុគមន៍ $f(x) = x^3 - 3x + 2$ មានក្រាបតំណាង $(C)$ ក្នុងប្លង់កូអរដោណេ $(O, \vec{i}, \vec{j})$ ។
			\begin{khmerenum}
				\item សិក្សាអថេរភាព និងសង់តារាងអថេរភាពនៃអនុគមន៍ $f$ ។
				\item រកកូអរដោណេចំណុចរបត់ $I$ នៃក្រាប $(C)$ ។
				\item សរសេរសមីការបន្ទាត់ប៉ះ $(T)$ ទៅនឹងក្រាប $(C)$ ត្រង់ចំណុច $I$ ។
				\item សង់ក្រាប $(C)$ និងបន្ទាត់ប៉ះ $(T)$ ក្នុងប្រព័ន្ធកូអរដោណេតែមួយ។
			\end{khmerenum}
		\end{circleenum}
	\end{sectionbox}

	\begin{notebox}
		\textbf{បម្រាម៖} បេក្ខជនមិនត្រូវបានអនុញ្ញាតឱ្យប្រើប្រាស់ម៉ាស៊ីនគិតលេខ ឬឯកសារផ្សេងៗចូលក្នុងបន្ទប់ប្រឡងឡើយ។
	\end{notebox}

\end{document}
